\setcounter{section}{13}

  \zwischenueberschrift{Soal latihan}

\definitionswort {Produk Kronecker}{}
dari matriks-matriks

\mavergleichskettedisp
{\vergleichskette
{ A
}
{ =} { \left(  a_{ij}  \right)_{1 \leq i \leq m,\, 1 \leq j \leq n}
}
{ } {
}
{ } {
}
{ } {
}
}
{}{}{}
dan

\mavergleichskettedisp
{\vergleichskette
{ B
}
{ =} { \left(  b_{k \ell }  \right)_{1 \leq k \leq p,\, 1 \leq \ell \leq r}
}
{ } {
}
{ } {
}
{ } {
}
}
{}{}{}
diberikan oleh

\mathdisp {\left(  a_{ij} \cdot b_{k \ell}  \right)_{1 \leq i \leq m,\,     1 \leq  k \leq p;\,   1 \leq j \leq n,     \, 1 \leq \ell \leq r   }} {  }

\inputaufgabegibtloesung
{}
{

Hitunglah
\definitionsverweis {produk Kronecker}{}{}
dari kedua matriks
\mathkor {} {\begin{pmatrix}  3   &   -4  \\  5  &  -2 \end{pmatrix}} {dan} {\begin{pmatrix}  -2  &   7   \\  6  &  3 \end{pmatrix}} {.}

}
{} {}

\inputaufgabe
{}
{

Misalkan $K$ sebuah
\definitionsverweis {lapangan}{}{}
dan

\mavergleichskettedisp
{\vergleichskette
{ A
}
{ =} { \left(  a_{ij}  \right)_{1 \leq i \leq m,\, 1 \leq j \leq n}
}
{ } {
}
{ } {
}
{ } {
}
}
{}{}{}
serta

\mavergleichskettedisp
{\vergleichskette
{ B
}
{ =} { \left(  b_{k \ell }  \right)_{1 \leq  k \leq p,\, 1 \leq \ell \leq r}
}
{ } {
}
{ } {
}
{ } {
}
}
{}{}{}
\definitionsverweis {matriks-matriks}{}{}
dengan
\definitionsverweis {pemetaan-pemetaan linear}{}{}
yang bersesuaian
\maabb {A} {K^n  } { K^m
} {}
dan
\maabb {B} {K^r } { K^p
} {.}
Tunjukkan bahwa
\definitionsverweis {produk tensor}{}{}
dari pemetaan-pemetaan linear tersebut, terhadap basis-basis

\mathbed {e_j \otimes e_\ell} {}
 {1 \leq j \leq n,\, 1 \leq \ell \leq r} {}
 {} {} {} {,}
dari
\mathl{K^n \otimes K^r}{} dan

\mathbed {e_i \otimes e_k} {}
 {1 \leq i \leq m,\, 1 \leq  k \leq p} {}
 {} {} {} {,}
dari
\mathl{K^m \otimes K^p}{,}
dideskripsikan oleh
\definitionsverweis {produk Kronecker}{}{}
dari $A$ dan $B$.

}
{} {}

\inputaufgabe
{}
{

Tunjukkan bahwa
\definitionsverweis {produk tensor}{}{}
dari
\definitionsverweis {pita Möbius}{}{}
dengan dirinya sendiri merupakan sebuah bundel garis trivial.

}
{} {}

\inputaufgabe
{}
{

Misalkan
\mathkor {} {M_1} {dan} {M_2} {}
\definitionsverweis {manifold-manifold berorientasi}{}{.}
Tunjukkan bahwa
\definitionsverweis {produk}{}{}

\mathl{M_1 \times M_2}{} merupakan sebuah manifold berorientasi
\zusatzklammer {dengan orientasi yang bergantung pada urutan pada \mathlk{\{1,2\}}{}} {} {.}

}
{} {}

\inputaufgabe
{}
{

Tunjukkan bahwa sfera berdimensi $1$, yaitu $S^1$, merupakan sebuah
\definitionsverweis {manifold diferensiabel}{}{}
yang
\definitionsverweis {dapat diorientasikan}{}{.}

}
{} {}

Misalkan
\mathkor {} {M} {dan} {N} {}
\definitionsverweis {manifold-manifold berorientasi}{}{}
dan
\maabbdisp {\varphi} { M } { N
} {}
sebuah
\definitionsverweis {pemetaan diferensiabel}{}{.}
Pemetaan $\varphi$ disebut \definitionswort {mempertahankan orientasi}{,} jika untuk setiap titik

\mavergleichskette
{\vergleichskette
{ P
}
{ \in }{ M
}
{  }{
}
{    }{
}
{   }{
}
}
{}{}{}
\definitionsverweis {pemetaan singgung}{}{}
\maabbdisp {T\varphi} { T_PM } { T_{\varphi(P)} N
} {}
bijektif dan
\definitionsverweis {mempertahankan orientasi}{}{.}

\inputaufgabe
{}
{

Tunjukkan bahwa
\definitionsverweis {pemetaan antipodal}{}{}
\maabbeledisp {\varphi} {S^1} {S^1
} { P } { -P
} {,}
\definitionsverweis {mempertahankan orientasi}{}{.}

}
{} {}

\inputaufgabe
{}
{

Misalkan $M$ sebuah
\definitionsverweis {manifold berorientasi}{}{.}
Tunjukkan bahwa pemetaan pertukaran
\maabbeledisp {\varphi} {M \times M} {M \times M
} {(P,Q)} {(Q,P)
} {,}
terhadap masing-masing
\definitionsverweis {orientasi produk}{}{,}
tidak harus
\definitionsverweis {mempertahankan orientasi}{}{.}

}
{} {}

\inputaufgabe
{}
{

Misalkan $X$ sebuah
\definitionsverweis {ruang topologis}{}{}
yang hanya terdiri atas
\definitionsverweis {berhingga banyak}{}{}
elemen. Tunjukkan bahwa $X$
\definitionsverweis {kompak}{}{.}

}
{} {}

\inputaufgabe
{}
{

Misalkan $X$ sebuah
\definitionsverweis {ruang topologis}{}{}
dan misalkan

\mavergleichskette
{\vergleichskette
{Y_1 , \ldots , Y_n
}
{ \subseteq }{ X
}
{  }{
}
{    }{
}
{   }{
}
}
{}{}{}
\definitionsverweis {himpunan-himpunan bagian kompak}{}{.}
Tunjukkan bahwa
\definitionsverweis {gabungan}{}{}

\mavergleichskette
{\vergleichskette
{ Y
}
{ = }{ \bigcup_{i = 1}^n Y_i
}
{  }{
}
{    }{
}
{   }{
}
}
{}{}{}
juga kompak.

}
{} {}

\inputaufgabegibtloesung
{}
{

Misalkan $X$ sebuah
\definitionsverweis {ruang kompak}{}{}
dan misalkan

\mavergleichskette
{\vergleichskette
{Y
}
{ \subseteq }{ X
}
{  }{
}
{    }{
}
{   }{
}
}
{}{}{}
sebuah
\definitionsverweis {himpunan bagian tertutup}{}{,}
yang dilengkapi dengan
\definitionsverweis {topologi terinduksi}{}{.}
Tunjukkan bahwa $Y$ juga kompak.

}
{} {}

\inputaufgabegibtloesung
{}
{

Misalkan $X$ sebuah
\definitionsverweis {ruang topologis}{}{}
\definitionsverweis {kompak}{}{}
yang tak kosong dan misalkan
\maabbdisp {f} { X } {\R
} {}
sebuah
\definitionsverweis {fungsi kontinu}{}{.}
Tunjukkan bahwa terdapat suatu

\mavergleichskette
{\vergleichskette
{ x
}
{ \in }{ X
}
{  }{
}
{    }{
}
{   }{
}
}
{}{}{}
dengan

\mathdisp {f(x) \geq f(x')  \text { untuk semua } x' \in X} {  }
.

}
{} {}

\inputaufgabe
{}
{

Misalkan
\maabbdisp {f} {\R} {S^1
} {}
sebuah
\definitionsverweis {pemetaan kontinu}{}{.}
Tunjukkan bahwa
\definitionsverweis {citra}{}{}
dari $f$
\definitionsverweis {homeomorfik}{}{}
dengan suatu
\definitionsverweis {interval}{}{}
terbuka, setengah terbuka, atau tertutup, atau dengan $S^1$.

}
{} {}

\inputaufgabe
{}
{

Misalkan
\maabbdisp {f} {S^1} {\R
} {}
sebuah
\definitionsverweis {pemetaan kontinu}{}{.}
Tunjukkan bahwa
\definitionsverweis {citra}{}{}
dari $f$
\definitionsverweis {homeomorfik}{}{}
dengan suatu
\definitionsverweis {interval tertutup}{}{.}

}
{} {}

\inputaufgabe
{}
{

Tunjukkan bahwa himpunan
\definitionsverweis {bilangan-bilangan riil}{}{} $\R$ tidak
\definitionsverweis {kompak terhadap tutupan}{}{.}

}
{} {}

\inputaufgabe
{}
{

Kita meninjau
\definitionsverweis {bilangan-bilangan asli}{}{}
$\N$ dan melengkapinya dengan
\definitionsverweis {metrik diskret}{}{.}
Tunjukkan bahwa $\N$
\definitionsverweis {
tertutup}{}{}
dan
\definitionsverweis {terbatas}{}{,}
tetapi tidak
\definitionsverweis {
kompak terhadap tutupan}{}{.}

}
{} {}

\inputaufgabegibtloesung
{}
{

Misalkan $X$ sebuah
\definitionsverweis {ruang metrik}{}{}
\definitionsverweis {kompak}{}{.}
Tunjukkan bahwa $X$
\definitionsverweis {
lengkap}{}{.}

}
{} {}

\inputaufgabe
{}
{

Tunjukkan bahwa terdapat sebuah
\definitionsverweis {keluarga terhitung}{}{}
dari
\definitionsverweis {bola-bola terbuka}{}{}
di $\R^n$ yang membentuk sebuah
\definitionsverweis {basis topologi}{}{.}

}
{} {}

\inputaufgabegibtloesung
{}
{

Tunjukkan bahwa himpunan

\mavergleichskettedisp
{\vergleichskette
{ M
}
{ =} { \left\{ (x,y,z) \in \R^3  \mid   x^2+y^4+z^6 = 1        \right\}
}
{ } {
}
{ } {
}
{ } {
}
}
{}{}{}
merupakan sebuah manifold diferensiabel kompak berdimensi dua.

}
{} {}

\inputaufgabegibtloesung
{}
{

Misalkan $M$ sebuah manifold topologis kompak berdimensi $d$,

\mavergleichskette
{\vergleichskette
{ d
}
{ \geq }{ 1
}
{  }{
}
{    }{
}
{   }{
}
}
{}{}{.}
Tunjukkan bahwa terdapat sebuah himpunan bagian terbuka terbatas

\mavergleichskette
{\vergleichskette
{ U
}
{ \subseteq }{ \R^d
}
{  }{
}
{    }{
}
{   }{
}
}
{}{}{}
dan sebuah pemetaan kontinu surjektif
\maabbdisp {\varphi} { U } { M
} {}.

}
{} {}

  \zwischenueberschrift{Soal untuk dikumpulkan}

\inputaufgabe
{4}
{

Tunjukkan bahwa sfera berdimensi $2$, yaitu $S^2$, merupakan sebuah
\definitionsverweis {manifold diferensiabel}{}{}
yang
\definitionsverweis {dapat diorientasikan}{}{.}

}
{} {}

\inputaufgabegibtloesung
{4}
{

Tunjukkan bahwa
\definitionsverweis {pemetaan antipodal}{}{}
\maabbeledisp {} {S^2} {S^2
} {(x,y,z)} {(-x,-y,-z)
} {,}
tidak
\definitionsverweis {mempertahankan orientasi}{}{.}

}
{} {}

\inputaufgabegibtloesung
{4}
{

Misalkan $X$ sebuah
\definitionsverweis {ruang Hausdorff}{}{}
dan misalkan

\mavergleichskette
{\vergleichskette
{Y
}
{ \subseteq }{ X
}
{  }{
}
{    }{
}
{   }{
}
}
{}{}{}
sebuah himpunan bagian yang dilengkapi dengan
\definitionsverweis {topologi terinduksi}{}{.}
Misalkan $Y$
\definitionsverweis {kompak}{}{.}
Tunjukkan bahwa $Y$
\definitionsverweis {tertutup}{}{}
di $X$.

}
{} {}

\inputaufgabe
{4}
{

Misalkan
\mathkor {} {X} {dan} {Y} {}
\definitionsverweis {ruang-ruang topologis}{}{}
dan misalkan
\maabbdisp {\varphi} { X } { Y
} {}
sebuah
\definitionsverweis {pemetaan kontinu}{}{.}
Misalkan $X$
\definitionsverweis {kompak}{}{.}
Tunjukkan bahwa
\definitionsverweis {citra}{}{}

\mavergleichskette
{\vergleichskette
{ \varphi(X)
}
{ \subseteq }{ Y
}
{  }{
}
{    }{
}
{   }{
}
}
{}{}{}
juga kompak.

}
{} {}

\inputaufgabe
{4}
{

Misalkan $V$ sebuah
\definitionsverweis {ruang vektor Euklides}{}{}
ber-
\definitionsverweis {dimensi}{}{}
$n$ dan misalkan
\definitionsverweis {produk}{}{}

\mavergleichskette
{\vergleichskette
{ V^n
}
{ = }{ V \times \cdots \times V
}
{  }{
}
{    }{
}
{   }{
}
}
{}{}{}
dilengkapi dengan
\definitionsverweis {topologi produk}{}{.}
Misalkan $I$ sebuah
\definitionsverweis {interval riil}{}{}
dan
\maabbdisp {\varphi} { I } { V^n
} {}
sebuah
\definitionsverweis {pemetaan kontinu}{}{}
dengan sifat bahwa

\mavergleichskettedisp
{\vergleichskette
{ \varphi(t)
}
{ =} { (\varphi_1(t), \varphi_2(t) , \ldots , \varphi_n(t))
}
{ } {
}
{ } {
}
{ } {
}
}
{}{}{}
untuk setiap

\mavergleichskette
{\vergleichskette
{ t
}
{ \in }{ I
}
{  }{
}
{    }{
}
{   }{
}
}
{}{}{}
merupakan sebuah
\definitionsverweis {basis}{}{}
dari $V$. Tunjukkan bahwa semua basis

\mathbed {\varphi(t)} {}
 {t \in I} {}
 {} {} {} {,}
merepresentasikan
\definitionsverweis {orientasi}{}{}
yang sama pada $V$.

}
{} {}
